% ============================================================
% §3 · Requisitos del usuario — Manual de Usuario K-QGMIP
% ============================================================

\section{Requisitos del usuario}
\label{sec:requisitos}

Esta sección lista los prerrequisitos de software, hardware y conocimiento necesarios para
ejecutar K-QGMIP. Cumplirlos garantiza resultados correctos y tiempos de ejecución razonables.

% ------------------------------------------------------------
\subsection{Requisitos de software}
% ------------------------------------------------------------

\begin{itemize}
  \item \textbf{Python 3.12 o superior.} El proyecto declara
        \texttt{requires-python = ">=3.12"} en su configuración. Versiones anteriores no
        están soportadas.

  \item \textbf{Gestor de dependencias \texttt{uv}.} Es el gestor recomendado por el proyecto
        para crear el entorno virtual e instalar todas las bibliotecas en una sola operación
        (ver Sección de Instalación y Configuración). Como alternativa puede usarse \texttt{pip}
        con el archivo \texttt{pyproject.toml}.

  \item \textbf{Dependencias declaradas en \texttt{pyproject.toml}.} La tabla siguiente
        resume las bibliotecas requeridas:
\end{itemize}

\begin{table}[h!]
\centering
\begin{tabular}{|l|l|p{7cm}|}
\hline
\textbf{Paquete} & \textbf{Versión mínima} & \textbf{Uso en K-QGMIP} \\
\hline
\texttt{numpy}        & $\geq$ 2.0.2  & Operaciones matriciales, TPM, N-Cubos \\
\texttt{scipy}        & $\geq$ 1.17.0 & Earth Mover's Distance (EMD) entre distribuciones \\
\texttt{pandas}       & $\geq$ 2.3.3  & Lectura de Excel y escritura de resultados en CSV \\
\texttt{openpyxl}     & $\geq$ 3.1.5  & Motor de lectura/escritura de archivos Excel \\
\texttt{pyphi}        & $\geq$ 1.2.0  & Referencia de validación para casos pequeños ($k=2$) \\
\texttt{pyinstrument} & $\geq$ 5.1.2  & Generación de reportes de rendimiento (HTML) \\
\texttt{colorama}     & $\geq$ 0.4.6  & Visualización colorizada de resultados en consola \\
\texttt{pyttsx3}      & $\geq$ 2.99   & Síntesis de voz al encontrar la solución óptima \\
\hline
\end{tabular}
\caption{Dependencias externas requeridas por K-QGMIP.}
\label{tab:dependencias}
\end{table}

% ------------------------------------------------------------
\subsection{Requisitos de hardware}
% ------------------------------------------------------------

El costo computacional de K-QGMIP escala con dos factores independientes: el tamaño del
sistema ($n$) y el número de particiones ($k$). La Tabla~\ref{tab:hardware} orienta al
usuario sobre los recursos mínimos recomendados según el rango de $n$; el impacto de $k$
se discute en la nota al pie de la tabla.

\begin{table}[h!]
\centering
\begin{tabular}{|c|c|p{7.5cm}|}
\hline
\textbf{Tamaño del sistema} & \textbf{RAM mínima} & \textbf{Perfil de máquina recomendado} \\
\hline
$n \leq 6$    & 4 GB  & Cualquier computador moderno (portátil académico básico) \\
$n = 8$--$10$ & 4 GB o superior & Portátil o escritorio estándar de uso académico \\
$n = 15$--$17$ & 16 GB & PC de escritorio con procesador de al menos 4 núcleos \\
$n = 20$--$22$ & 32 GB & PC de alto rendimiento \\
$n \geq 25$   &  32 GB o superior & PC de alto rendimiento \\
\hline
\end{tabular}
\caption{Recursos de hardware recomendados según tamaño del sistema analizado ($k=2$).}
\label{tab:hardware}
\end{table}

\textbf{Efecto de $k$ sobre el costo computacional:} los recursos de la
Tabla~\ref{tab:hardware} aplican al caso base $k=2$. Al incrementar $k$, lo
que puede aumentar considerablemente el tiempo de ejecución y el uso de memoria para un
mismo $n$. Como regla general: a mayor $k$, mayor costo; se recomienda comenzar siempre
con $k=2$ para estimar la carga computacional del sistema antes de escalar a valores
superiores.

\textbf{GPU}: no es un requisito obligatorio. K-QGMIP corre íntegramente en CPU. Sin embargo,
si el entorno dispone de múltiples núcleos, la paralelización de evaluaciones de particiones
puede configurarse para aprovecharlos y reducir los tiempos indicados.

% ------------------------------------------------------------
\subsection{Requisitos de conocimiento}
% ------------------------------------------------------------

Para utilizar K-QGMIP no se exige experiencia en algoritmos avanzados ni en matemáticas de
IIT 4.0, pero sí se asume la siguiente familiaridad básica:

\begin{itemize}
  \item \textbf{Python básico}: capacidad de ejecutar un script (\texttt{python main.py}),
        instalar paquetes y leer mensajes de error de la consola.
  \item \textbf{Línea de comandos o IDE}: saber navegar directorios y abrir un terminal
        (símbolo del sistema, PowerShell o terminal de VS Code / PyCharm).
  \item \textbf{Formato de datos}: entender qué es un archivo CSV y un archivo Excel
        (\texttt{.xlsx}), puesto que la TPM se carga desde CSV y los casos de prueba
        en lote se definen en Excel.
  \item \textbf{Nociones de IIT} (\textit{opcional pero recomendado}): familiaridad con
        los conceptos de matriz de transición probabilística, estado inicial del sistema
        y la idea general de información integrada $\varphi$, para poder interpretar
        correctamente los resultados obtenidos.
\end{itemize}
